\documentclass[11pt]{article}
\usepackage[a4paper,margin=22mm]{geometry}
\usepackage{fontspec}
\setmainfont{DejaVu Serif}
\setsansfont{DejaVu Sans}
\usepackage{microtype}
\usepackage{xcolor}
\usepackage{hyperref}
\usepackage{enumitem}
\usepackage{parskip}
\definecolor{Accent}{HTML}{971D32}
\definecolor{Muted}{HTML}{666666}
\hypersetup{colorlinks=true,urlcolor=Accent,linkcolor=Accent}
\setlist{leftmargin=1.6em,itemsep=0.3em,topsep=0.35em}
\setcounter{tocdepth}{2}
\renewcommand{\contentsname}{Contents}
\newcommand{\sectionrule}{\par\vspace{-0.5em}\noindent\textcolor{Accent}{\rule{\linewidth}{0.6pt}}\par}
\begin{document}

\begin{center}
{\sffamily\bfseries\color{Accent} TOEFL WORD OF THE DAY\par}
\vspace{8mm}
{\fontsize{42}{48}\selectfont\bfseries propriety\par}
\vspace{2mm}
{\Large\sffamily\color{Accent} /prəˈpraɪəti/\par}
\vspace{2mm}
{\small\sffamily Local audio phrase: propriety\par}
{\small\sffamily Audio file: audios/propriety.mp3\par}
\vspace{2mm}
{\large\sffamily\color{Muted} uncountable noun; plural countable noun\par}
\vspace{5mm}
{\sffamily correct and appropriate behavior \textbullet\ accepted rules of polite behavior\par}
\vspace{6mm}
{\large Propriety refers to behavior or decisions that meet accepted moral, social, professional, or procedural standards.\par}
\vspace{6mm}
\begin{minipage}{0.84\textwidth}
\itshape “The committee questioned the propriety of allowing the director to approve a contract involving his own company.”
\end{minipage}\par
\vspace{10mm}
\textbf{Date:} September 1st 2026\quad\textbullet\quad \textbf{Level:} B1--mid-B2 TOEFL
\vfill
Created and compiled by \textbf{Amir Kooshky}\par
\vspace{2mm}
\href{https://t.me/KooshkyTOEFL}{Telegram Channel}\quad\textbullet\quad
\href{https://www.instagram.com/kooshkytoefl}{Instagram}\quad\textbullet\quad
\href{https://t.me/KooshkyTOEFL\_pv}{Personal Telegram}
\end{center}
\newpage
\tableofcontents
\newpage

\section{Meanings and usage}
\sectionrule
\subsection{Meaning 1: Correct and appropriate behavior}
\textbf{Part of speech:} uncountable noun

\textbf{IPA:} /prəˈpraɪəti/

\textbf{Pronunciation phrase:} propriety

\textbf{Local audio file:} audios/propriety.mp3

\textbf{Register:} formal; common in legal, ethical, professional, political, and institutional contexts

\textbf{Definition:} behavior that is considered morally, socially, or professionally correct and appropriate

\textbf{In simpler words:} being correct and appropriate

\textbf{Typical pattern:} the propriety of + action or decision

\subsubsection{Natural examples}
\begin{enumerate}
  \item The committee questioned the propriety of allowing the director to approve a contract involving his own company.
  \item Employees are expected to observe professional propriety when dealing with confidential information.
  \item The debate focused on the propriety of using public funds for the project.
  \item Several scholars raised concerns about the ethical propriety of conducting the experiment without full consent.
  \item The judge considered whether the official's actions met accepted standards of propriety.
  \item Critics questioned the propriety of giving government contracts to close political allies.
  \item The organization introduced new rules to ensure propriety in its financial decisions.
  \item Questions about the propriety of the procedure delayed the final decision.
\end{enumerate}

\subsubsection{Nuance and implication}
\textbf{Central nuance:} Propriety is about whether behavior or a decision fits accepted standards of what is proper, appropriate, ethical, or respectable.

\textbf{Usual implication:} The word often appears when people are evaluating whether an action crosses a moral, professional, or procedural boundary.

\textbf{Compared with impropriety:} Propriety means behavior that meets accepted standards, while impropriety means behavior that violates those standards.

\subsubsection{Grammar notes}
\textbf{How to use it:} In this meaning, propriety is normally uncountable. It often appears in formal expressions such as the propriety of an action or questions about propriety.

\textbf{What can follow it:} It is commonly followed by of and then the action, decision, or behavior being evaluated.

\textbf{Common preposition:} Use the propriety of + noun or -ing form, as in the propriety of the decision or the propriety of accepting the gift.

\textbf{Noun use:} Do not normally use a or an before propriety in this meaning.

\textbf{Position in a sentence:} It often appears after words such as question, doubt, consider, discuss, ensure, or observe.

\subsubsection{Useful patterns}
\textbf{Pattern:} the propriety of + noun or -ing form

\textbf{Use:} Use this when evaluating whether an action or decision is appropriate.

\textbf{Example:} The board debated the propriety of accepting donations from the company.

\textbf{Pattern:} question the propriety of + action

\textbf{Use:} Use this when expressing doubt about whether something is proper or acceptable.

\textbf{Example:} Several members questioned the propriety of holding the vote in private.

\textbf{Pattern:} standards of propriety

\textbf{Use:} Use this for accepted rules about proper conduct.

\textbf{Example:} Public officials are expected to follow high standards of propriety.

\textbf{Pattern:} professional propriety

\textbf{Use:} Use this for behavior that meets accepted professional standards.

\textbf{Example:} The investigation examined whether the lawyer had maintained professional propriety.

\textbf{Pattern:} ethical propriety

\textbf{Use:} Use this when judging whether an action is ethically acceptable.

\textbf{Example:} The researchers debated the ethical propriety of the proposed experiment.

\subsubsection{Common wrong usage}
\paragraph{Correction 1}
\textbf{Wrong:} The committee discussed about the propriety of the decision.

\textbf{Correct:} The committee discussed the propriety of the decision.

\textbf{Why:} Discuss normally takes its object directly, without about.

\paragraph{Correction 2}
\textbf{Wrong:} They questioned a propriety of using public funds.

\textbf{Correct:} They questioned the propriety of using public funds.

\textbf{Why:} In this meaning, propriety is normally uncountable and commonly appears as the propriety of.

\paragraph{Correction 3}
\textbf{Wrong:} Propriety means that something is legal.

\textbf{Correct:} Propriety means that something is considered appropriate or proper.

\textbf{Why:} An action can be legal but still be considered improper. Propriety is broader than legality.

\paragraph{Correction 4}
\textbf{Wrong:} The official acted with many proprieties.

\textbf{Correct:} The official acted with propriety.

\textbf{Why:} When referring generally to proper conduct, use the uncountable noun propriety.

\subsection{Meaning 2: Accepted rules of polite behavior}
\textbf{Part of speech:} plural countable noun

\textbf{IPA:} /prəˈpraɪəti/

\textbf{Pronunciation phrase:} propriety

\textbf{Local audio file:} audios/propriety.mp3

\textbf{Register:} formal and somewhat old-fashioned; mainly used in discussions of etiquette, social customs, literature, and history

\textbf{Definition:} the accepted rules of polite and socially appropriate behavior

\textbf{In simpler words:} rules of proper social behavior

\textbf{Typical pattern:} observe, follow, ignore, or respect + the proprieties

\subsubsection{Natural examples}
\begin{enumerate}
  \item Guests were expected to observe the proprieties of formal diplomatic events.
  \item He cared little about social proprieties and often spoke very directly.
  \item The novel shows how strictly the proprieties of upper-class society controlled personal behavior.
  \item She followed all the usual proprieties when meeting the ambassador.
  \item Changes in social attitudes have weakened some of the proprieties that governed public behavior in the past.
  \item The ceremony was conducted according to the traditional proprieties of the community.
\end{enumerate}

\subsubsection{Nuance and implication}
\textbf{Central nuance:} The plural proprieties refers to specific conventions about how people are expected to behave politely or respectably.

\textbf{Usual implication:} These rules are usually created by social tradition rather than formal law.

\textbf{Compared with breach of etiquette:} Observing the proprieties means following accepted social conventions, while a breach of etiquette means violating one of them.

\subsubsection{Grammar notes}
\textbf{How to use it:} This meaning is normally plural and commonly appears as the proprieties or social proprieties.

\textbf{What can follow it:} It often appears after verbs such as observe, follow, respect, maintain, or ignore.

\textbf{Noun use:} Use the plural proprieties when referring to particular social conventions or rules of polite behavior.

\textbf{Position in a sentence:} The expression observe the proprieties is especially common in formal writing.

\subsubsection{Useful patterns}
\textbf{Pattern:} observe the proprieties

\textbf{Use:} Use this when someone follows accepted rules of polite or respectable behavior.

\textbf{Example:} Despite their disagreement, both diplomats carefully observed the proprieties.

\textbf{Pattern:} social proprieties

\textbf{Use:} Use this for accepted social conventions about correct behavior.

\textbf{Example:} The novel frequently challenges the social proprieties of the period.

\textbf{Pattern:} ignore the proprieties

\textbf{Use:} Use this when someone does not follow expected social conventions.

\textbf{Example:} He ignored the proprieties and openly criticized his host during dinner.

\subsubsection{Common wrong usage}
\paragraph{Correction 1}
\textbf{Wrong:} She followed the propriety of formal dinners.

\textbf{Correct:} She followed the proprieties of formal dinners.

\textbf{Why:} When referring to specific rules of social behavior, the plural proprieties is normally used.

\paragraph{Correction 2}
\textbf{Wrong:} Social proprieties are official laws.

\textbf{Correct:} Social proprieties are accepted conventions of proper behavior.

\textbf{Why:} Proprieties are social expectations, not necessarily legal rules.

\section{Word family}
\sectionrule
\subsection{proper}
\textbf{Part of speech:} adjective

\textbf{IPA:} /ˈprɑːpər/

\textbf{Definition:} correct, suitable, or appropriate for a particular situation

\textbf{Relationship to the headword:} Proper expresses the basic idea of being correct or appropriate that is contained in propriety.

\textbf{Grammar pattern:} proper + behavior, procedure, method, care, or use

\subsubsection{Examples}
\begin{enumerate}
  \item Researchers must follow proper procedures when collecting personal data.
  \item The equipment cannot function safely without proper maintenance.
\end{enumerate}

\subsection{impropriety}
\textbf{Part of speech:} countable or uncountable noun

\textbf{IPA:} /ˌɪmprəˈpraɪəti/

\textbf{Definition:} behavior that is considered improper, unethical, or unacceptable

\textbf{Relationship to the headword:} Impropriety is the opposite form of propriety and refers to behavior that violates accepted standards.

\textbf{Grammar pattern:} financial, ethical, professional, or procedural + impropriety

\textbf{Usage note:} The plural improprieties is possible when referring to separate improper acts.

\subsubsection{Examples}
\begin{enumerate}
  \item The investigation found no evidence of financial impropriety.
  \item The official resigned after allegations of serious impropriety.
  \item Several procedural improprieties were identified during the review.
\end{enumerate}

\section{Common collocations}
\sectionrule
\subsection{the propriety of}
\textbf{Applies to:} Correct and appropriate behavior

\textbf{Meaning:} whether an action or decision is proper or appropriate

\textbf{Pattern:} the propriety of + noun or -ing form

\textbf{Example:} The committee debated the propriety of accepting the donation.

\subsection{question the propriety of}
\textbf{Applies to:} Correct and appropriate behavior

\textbf{Meaning:} express doubt about whether something is appropriate

\textbf{Example:} Critics questioned the propriety of allowing officials to regulate companies in which they owned shares.

\subsection{standards of propriety}
\textbf{Applies to:} Correct and appropriate behavior

\textbf{Meaning:} accepted standards of proper behavior

\textbf{Example:} Judges are expected to maintain strict standards of propriety.

\subsection{professional propriety}
\textbf{Applies to:} Correct and appropriate behavior

\textbf{Meaning:} behavior that meets accepted professional standards

\textbf{Example:} The guidelines are intended to protect professional propriety and public trust.

\subsection{ethical propriety}
\textbf{Applies to:} Correct and appropriate behavior

\textbf{Meaning:} the ethical acceptability of an action

\textbf{Example:} The study raised questions about the ethical propriety of experimenting on vulnerable populations.

\subsection{observe the proprieties}
\textbf{Applies to:} Accepted rules of polite behavior

\textbf{Meaning:} follow accepted rules of polite or respectable behavior

\textbf{Example:} Both sides observed the proprieties during the formal negotiations.

\subsection{social proprieties}
\textbf{Applies to:} Accepted rules of polite behavior

\textbf{Meaning:} accepted conventions governing social behavior

\textbf{Example:} The author criticized the rigid social proprieties of the nineteenth century.

\subsection{financial impropriety}
\textbf{Applies to:} Correct and appropriate behavior

\textbf{Meaning:} dishonest or improper behavior involving money

\textbf{Example:} The audit found evidence of financial impropriety.

\textbf{Usage note:} Impropriety is the opposite word-family member.

\section{Synonyms and antonyms}
\sectionrule
\subsection{Close synonyms}
\subsubsection{appropriateness}
\textbf{Part of speech:} uncountable noun

\textbf{Applies to:} Correct and appropriate behavior

\textbf{Shared meaning:} Both can refer to whether something is suitable or proper.

\textbf{Difference:} Appropriateness is the general quality of being suitable for a situation. Propriety is more formal and often carries moral, professional, or social judgment.

\textbf{Example:} The committee considered the appropriateness of using the material in a classroom setting.

\subsubsection{decorum}
\textbf{Part of speech:} uncountable noun

\textbf{Applies to:} Accepted rules of polite behavior

\textbf{Shared meaning:} Both can refer to behavior that follows accepted social standards.

\textbf{Difference:} Decorum emphasizes polite, dignified behavior, especially in formal situations. Propriety is broader and can refer to social, ethical, or professional correctness.

\textbf{Example:} Members of the audience were expected to maintain decorum during the ceremony.

\subsubsection{correctness}
\textbf{Part of speech:} uncountable noun

\textbf{Applies to:} Correct and appropriate behavior

\textbf{Shared meaning:} Both can express the idea of meeting an accepted standard.

\textbf{Difference:} Correctness is broader and can refer to factual, grammatical, or procedural accuracy. Propriety specifically concerns whether conduct or a decision is proper or acceptable.

\textbf{Example:} The researchers checked the correctness of the calculations.

\subsection{Direct or useful antonyms}
\subsubsection{impropriety}
\textbf{Part of speech:} countable or uncountable noun

\textbf{Applies to:} Correct and appropriate behavior

\textbf{Opposition:} Propriety means proper or appropriate conduct, while impropriety means conduct that violates accepted moral, professional, or social standards.

\textbf{Example:} The investigation found evidence of serious financial impropriety.

\subsubsection{indecorum}
\textbf{Part of speech:} uncountable noun

\textbf{Applies to:} Accepted rules of polite behavior

\textbf{Opposition:} Indecorum means behavior that violates accepted standards of politeness or dignity, while observing the proprieties means following those standards.

\textbf{Usage note:} Indecorum is formal and relatively uncommon.

\textbf{Example:} His behavior at the formal ceremony was criticized as indecorum.

\section{Words learners may confuse}
\sectionrule
\subsection{property}
\textbf{Part of speech:} noun

\textbf{Why learners confuse them:} The words are spelled and pronounced somewhat similarly but have very different modern meanings.

\textbf{Key difference:} Propriety means proper or acceptable behavior. Property means something owned by a person or organization, or a characteristic of a substance or object.

\textbf{propriety:} The committee questioned the propriety of the arrangement.

\textbf{property:} The company owns property in several cities.

\subsection{proprietary}
\textbf{Part of speech:} noun versus adjective

\textbf{Why learners confuse them:} The words share a similar form and historical connection with ideas of what belongs or is proper.

\textbf{Key difference:} Propriety concerns proper or acceptable conduct. Proprietary means relating to ownership or to something legally owned and controlled by a person or company.

\textbf{propriety:} There were doubts about the propriety of the official's behavior.

\textbf{proprietary:} The company refused to reveal its proprietary technology.

\subsection{impropriety}
\textbf{Part of speech:} noun

\textbf{Why learners confuse them:} Impropriety is formed by adding the negative prefix im- to propriety.

\textbf{Key difference:} Propriety means appropriate or acceptable conduct. Impropriety means conduct that is inappropriate, unethical, or unacceptable.

\textbf{propriety:} The official acted with complete propriety.

\textbf{impropriety:} The audit uncovered evidence of financial impropriety.

\section{Etymology}
\sectionrule
\textbf{Origin:} Propriety comes through French from Latin proprietas, which referred to ownership, a particular character, or what properly belongs to something.

\textbf{Development:} The word gradually developed the sense of what is proper or appropriate, especially in behavior and social conduct.

\textbf{Memory link:} Think of propriety as behaving in the way that is considered proper.

\section{Practice exercises}
\sectionrule
\subsection{Word family choice}
Choose the best member of the word family for each blank.

\begin{enumerate}
  \item The investigation found no evidence of financial \_\_\_\_\_.
  \item Employees are expected to follow \_\_\_\_\_ procedures when handling confidential records.
  \item The committee questioned the \_\_\_\_\_ of allowing the official to participate in the decision.
\end{enumerate}

\subsection{Correct or incorrect?}
Decide whether each sentence is correct. Correct the inaccurate ones.

\begin{enumerate}
  \item The committee questioned the propriety of accepting the donation.
  \item The official's behavior showed a propriety.
  \item Guests were expected to observe the social proprieties.
  \item Propriety simply means that an action is legal.
  \item The audit raised concerns about financial impropriety.
\end{enumerate}

\subsection{Collocation practice}
Complete each sentence with a natural collocation.

\begin{enumerate}
  \item Several members questioned the propriety \_\_\_\_\_ holding the meeting in private.
  \item Public officials are expected to maintain high standards \_\_\_\_\_ propriety.
  \item The guests carefully observed the social \_\_\_\_\_.
  \item The audit uncovered evidence of financial \_\_\_\_\_.
\end{enumerate}

\subsection{Complete answer key}
\subsubsection{Word family choice}
\begin{enumerate}
  \item \textbf{impropriety.} The investigation found no evidence of financial impropriety. \emph{Impropriety means improper or unethical behavior.}
  \item \textbf{proper.} Employees are expected to follow proper procedures when handling confidential records. \emph{Use the adjective proper before procedures.}
  \item \textbf{propriety.} The committee questioned the propriety of allowing the official to participate in the decision. \emph{The propriety of means whether an action is proper or appropriate.}
\end{enumerate}

\subsubsection{Correct or incorrect?}
\begin{enumerate}
  \item \textbf{Correct} \emph{The propriety of is a standard formal pattern for evaluating whether an action is appropriate.}
  \item \textbf{Incorrect}. The official's behavior showed propriety. \emph{When propriety means proper conduct in general, it is normally uncountable.}
  \item \textbf{Correct} \emph{The plural proprieties can refer to accepted social conventions.}
  \item \textbf{Incorrect}. Propriety means that an action is considered proper or appropriate. \emph{Propriety can concern ethical, professional, social, or procedural standards, not just legality.}
  \item \textbf{Correct} \emph{Financial impropriety is a common formal expression for improper or unethical financial behavior.}
\end{enumerate}

\subsubsection{Collocation practice}
\begin{enumerate}
  \item \textbf{of.} Several members questioned the propriety of holding the meeting in private. \emph{Use the propriety of + noun or -ing form.}
  \item \textbf{of.} Public officials are expected to maintain high standards of propriety. \emph{Standards of propriety means accepted standards of proper conduct.}
  \item \textbf{proprieties.} The guests carefully observed the social proprieties. \emph{Social proprieties are accepted rules of polite behavior.}
  \item \textbf{impropriety.} The audit uncovered evidence of financial impropriety. \emph{Financial impropriety means improper or unethical behavior involving money.}
\end{enumerate}

\vfill
\begin{center}
\small\color{Muted} Created and compiled by \textbf{Amir Kooshky}\\
\href{https://t.me/KooshkyTOEFL}{Telegram Channel}\quad\textbullet\quad
\href{https://www.instagram.com/kooshkytoefl}{Instagram}\quad\textbullet\quad
\href{https://t.me/KooshkyTOEFL\_pv}{Personal Telegram}
\end{center}
\end{document}
